%% PAPER 3 — DRAFT v0.1
%% A Universal Quartic Dispersion Framework: From Kerr Spacetime
%% to Condensed-Matter Analogs
%%
%% Companion to: Paper 1 (Zenodo 10.5281/zenodo.22018715, CQG-117140)
%%               Paper 2 (Zenodo 10.5281/zenodo.22051427)
%% Status: DRAFT — internal consolidation, not for distribution

\documentclass[12pt,a4paper]{article}
\usepackage{amsmath,amssymb,amsthm}
\usepackage{graphicx,hyperref,geometry,natbib,booktabs,xcolor}
\geometry{margin=2.5cm}

\newtheorem{theorem}{Theorem}
\newtheorem{result}{Result}
\newtheorem{remark}{Remark}
\newcommand{\TODO}[1]{\textcolor{red}{\textbf{[TODO: #1]}}}
\newcommand{\Lam}{\Lambda}

\begin{document}

\title{%
  A Universal Quartic Dispersion Framework:\\
  Numerical Validation, Multi-Messenger Time-of-Flight Bounds,\\
  and Condensed-Matter Analog Mappings
}
\author{%
  Dimitar Kretski\\
  Center for Hydro- and Aerodynamics,\\
  Bulgarian Academy of Sciences, Varna, Bulgaria\\
  \texttt{ORCID: 0000-0001-5108-2243}
}
\date{\today\ --- \textbf{DRAFT v0.1}}
\maketitle

\begin{abstract}
We consolidate and numerically validate the one-parameter quartic
dispersion framework $\omega(k)=ck\sqrt{1+\Lambda k^2}$ introduced in
Papers 1--2, and extend it in three directions.
\textbf{(1) Numerical infrastructure:} we implement and validate a
spectral solver for the corresponding wave equation, verify the
discrete biharmonic operator $D_n^4$ converges to the continuum limit
at the expected $O(h^2)$ order, and demonstrate that $\Lambda$ can be
recovered from noisy $(k,\omega)$ data to sub-percent accuracy.
\textbf{(2) Corrected Fermi-LAT audit:} we identify and correct a
factor-of-$3/2$ normalization error in the group-velocity relation
used in Paper 1's Fermi-LAT bound, tightening it from
$\Lambda<2.16\times10^{-53}\,\mathrm{m^2}$ to the audited value
$\Lambda<1.44\times10^{-53}\,\mathrm{m^2}$.
\textbf{(3) Condensed-matter analog mappings:} we show that the same
dispersion form appears in Bogoliubov excitations of Bose--Einstein
condensates ($\Lambda=\xi^2/4$), an analytically verified mapping, and
we release open-source fitting tools so that this claim is directly
falsifiable by experimentalists. We also examined a proposed mapping to
trigonal-warping-corrected Dirac materials and one to photonic
$k$-space dispersion; on closer inspection neither currently rests on a
literature-supported derivation (Section~\ref{sec:analogs}), and we
report this honestly rather than presenting them as validated.
We emphasize throughout the distinction between an \emph{audited
observational bound} (Fermi-LAT), a \emph{validated theoretical
sensitivity ceiling} (EHT), and an \emph{illustrative forecast}
(gravitational-wave multi-band observations) --- only the first
constitutes a genuine fitted constraint from published data.
\end{abstract}

\newpage\tableofcontents\newpage

%% ════════════════════════════════════════════════════════════════════════════
\section{Introduction}
%% ════════════════════════════════════════════════════════════════════════════

Papers 1 and 2 established a one-parameter Hamiltonian dispersion
framework,
\begin{equation}
  H = \tfrac{1}{2}\bigl[g^{\mu\nu}k_\mu k_\nu + \Lambda\,k_{\rm loc}^4\bigr],
\end{equation}
and derived its consequences in Schwarzschild and Kerr spacetime: two
proved theorems on field/worldline non-equivalence (Paper 1), and four
results on photon-ring chromaticity, frequency-dependent shadow
sensitivity, superradiance invariance, and eikonal quasinormal modes
(Paper 2).

The present paper serves three purposes:

\begin{enumerate}
  \item \textbf{Close the numerical-infrastructure gap.} Papers 1--2
    used custom root-finding and implicit-differentiation code for each
    specific calculation (photon rings, QNM eikonal parameters). Here we
    build and validate a general-purpose PDE solver for the underlying
    wave equation, independent of any particular spacetime, and confirm
    it reproduces the known dispersion relation to controlled numerical
    accuracy.

  \item \textbf{Audit and correct the Fermi-LAT bound.} A normalization
    error in the group-velocity relation used in Paper 1 is identified
    and corrected here (Section~\ref{sec:fermi}).

  \item \textbf{Extend the physical scope.} We show that the flat-space
    limit of the dispersion relation, $\omega=ck\sqrt{1+\Lambda k^2}$,
    is mathematically identical to the established Bogoliubov dispersion
    law for BEC excitations, and we release open-source fitting tools so
    that this claim is directly falsifiable by experimentalists. We also
    audit two further proposed condensed-matter mappings (Dirac
    materials, photonic crystals) and report that neither currently
    survives scrutiny in its present form (Section~\ref{sec:analogs}).
\end{enumerate}

\begin{remark}
This paper does not claim a unified physical origin for $\Lambda$ across
gravitational and condensed-matter contexts. It claims only that the
\emph{same functional form} of dispersion relation appears in each
context, with different physical identifications of the parameter. A
single numerical value of $\Lambda$ that is simultaneously consistent
across all probes would be evidence for common origin; this has not been
established and is not claimed here.
\end{remark}

%% ════════════════════════════════════════════════════════════════════════════
\section{Numerical Wave-Equation Solver}
\label{sec:solver}
%% ════════════════════════════════════════════════════════════════════════════

\subsection{Flat-space wave equation}

The flat-space limit of the Hamiltonian (with ZAMO/comoving observer
reducing to the ordinary observer) gives the real-space PDE
\begin{equation}\label{eq:pde}
  \partial_t^2\phi = c^2\nabla^2\phi - c^2\Lambda\,\nabla^4\phi,
\end{equation}
i.e.\ the standard wave equation with a biharmonic correction. We
implement a 2D periodic-domain Fourier spectral solver
(\texttt{wave\_equation\_2D\_solver.py}) in which the spatial operator
is exact (machine precision) and the only numerical error comes from
finite time-sampling in the frequency-extraction step.

\subsection{Validation: recovered dispersion relation}

Table~\ref{tab:dispersion_test} shows the numerically extracted
$\omega(k)$ for a plane-wave initial condition, compared to the exact
formula, at several $\Lambda$ values (grid: $64\times64$, wavevector
$(k_x,k_y)=(3,4)$, so $k=5$).

\begin{table}[h]
\centering
\caption{Numerical vs.\ exact dispersion relation, $k=5$.}
\label{tab:dispersion_test}
\begin{tabular}{lccc}
\toprule
$\Lambda$ & $\omega_{\rm exact}$ & $\omega_{\rm numerical}$ & relative error \\
\midrule
$0.00$ & 5.00000000 & 4.98333333 & $3.33\times10^{-3}$ \\
$0.01$ & 5.59016994 & 5.57153604 & $3.33\times10^{-3}$ \\
$0.05$ & 7.50000000 & 7.47500000 & $3.33\times10^{-3}$ \\
$0.10$ & 9.35414347 & 9.32296299 & $3.33\times10^{-3}$ \\
\bottomrule
\end{tabular}
\end{table}

The residual $3.33\times10^{-3}$ relative error is uniform across
$\Lambda$ and is entirely attributable to FFT frequency-bin discretization
in the time domain (finite number of oscillation periods sampled); it is
not a spatial discretization error, since the spectral spatial operator
is exact. Section~\ref{sec:convergence} confirms this error scales as
$O(1/N_{\rm samples})$.

\subsection{Discrete biharmonic operator $D_n^4$}

For finite-difference (as opposed to spectral) implementations, the
biharmonic term in Eq.~\eqref{eq:pde} is approximated by
$D_n^4=(D_n^2)^2$, the composition of the standard 3-point discrete
Laplacian. We verify (\texttt{paper3\_Dn4\_test.py}):

\begin{itemize}
  \item[\checkmark] \textbf{Symbol convergence:} the discrete Fourier
    symbol $[-(2/h^2)(1-\cos kh)]^2$ converges to the continuum $k^4$ at
    measured order $2.000$ (5 grid refinements, $h\in[0.01,0.5]$),
    matching the expected order for a standard 5-point stencil.
  \item[\checkmark] \textbf{Eigenfunction check:} direct matrix
    application of $D_n^4$ to a discretized plane wave reproduces the
    discrete-symbol prediction to $2\times10^{-10}$ (machine precision),
    confirming the matrix construction is correct independent of the
    continuum-limit question.
\end{itemize}

\subsection{Grid and temporal convergence}
\label{sec:convergence}

Table~\ref{tab:convergence} confirms $O(1/N)$ convergence of the
frequency-extraction error with the number of time samples, consistent
with standard FFT binning behavior and ruling out spurious numerical
drift in the solver.

\begin{table}[h]
\centering
\caption{Temporal resolution convergence ($\Lambda=0.05$, fixed spatial
grid).}
\label{tab:convergence}
\begin{tabular}{lcc}
\toprule
$N_{\rm samples}$ & relative error & convergence order \\
\midrule
100    & $1.00\times10^{-2}$ & --    \\
300    & $3.33\times10^{-3}$ & $-1.000$ \\
1000   & $1.00\times10^{-3}$ & $-1.000$ \\
3000   & $3.33\times10^{-4}$ & $-1.000$ \\
10000  & $1.00\times10^{-4}$ & $-1.000$ \\
\bottomrule
\end{tabular}
\end{table}

\subsection{$\Lambda$ recovery from noisy data}

We further test the inverse problem directly relevant to
experimentalists: given noisy $(k,\omega)$ pairs, can $\Lambda$ be
recovered by least-squares fitting? Using synthetic data with
$\Lambda_{\rm true}=0.07$ and $0.1\%$ Gaussian noise on $\omega$, the
fit recovers $\Lambda_{\rm fit}=0.069771$, a relative error of $0.33\%$
--- well within the noise level, confirming the fitting procedure
introduces no significant bias.

%% ════════════════════════════════════════════════════════════════════════════
\section{Corrected Fermi-LAT Bound}
\label{sec:fermi}
%% ════════════════════════════════════════════════════════════════════════════

\subsection{The normalization error}

Paper 1 derived a bound on $\Lambda$ from Fermi-LAT gamma-ray-burst
timing \citep{Abdo2009} by equating the fractional group-velocity
deviation $(v_g-c)/c$ to the standard linear-in-$\Lambda$ term of the
dispersion relation. Auditing this derivation reveals that the correct
relation is
\begin{equation}
  \frac{v_g - c}{c} = \frac{3}{2}\bigl(D-1\bigr),
\end{equation}
where $D=v_g/c$ was (incorrectly) treated as directly proportional to
$\Lambda\omega^2$ without the leading numerical factor $3/2$ that arises
from differentiating $\omega=ck\sqrt{1+\Lambda k^2}$ to obtain the group
velocity (Paper 1, Eq.~for $v_g$; the factor originates from
$d/dk[k\sqrt{1+\Lambda k^2}]$ at leading order in $\Lambda k^2$).

\subsection{Corrected bound}

\begin{result}[Audited Fermi-LAT bound]
Correcting for the $3/2$ factor gives
\begin{equation}
  \Lambda_{\rm corrected} = \frac{2}{3}\,\Lambda_{\rm original},
\end{equation}
i.e.\ a $33.3\%$ tightening of the bound:
\begin{equation}
  \boxed{\Lambda_{\rm Fermi} = 1.4421\times10^{-53}\ {\rm m^2}}
\end{equation}
in place of the previously quoted $2.1632\times10^{-53}\,{\rm m^2}$. The
order of magnitude, $\Lambda_{\rm Fermi}=O(10^{-53}\,{\rm m^2})$, is
unchanged.
\end{result}

This correction does not affect any result in Paper 2 (Kerr photon
rings, chromatic shadow, superradiance, eikonal QNMs), which do not
depend on the Fermi-LAT normalization.

%% ════════════════════════════════════════════════════════════════════════════
\section{Consolidated Multi-Probe Summary}
\label{sec:consolidated}
%% ════════════════════════════════════════════════════════════════════════════

Table~\ref{tab:probes} consolidates the three observational/theoretical
probes of $\Lambda$ developed across Papers 1--3, with an explicit
statement of evidentiary status for each.

\begin{table}[h]
\centering
\caption{Consolidated $\Lambda$ probes. See remarks below for status
definitions.}
\label{tab:probes}
\begin{tabular}{lllcl}
\toprule
Probe & Observable & Scaling & Result & Status \\
\midrule
Fermi-LAT     & time-of-flight            & $\Lambda\omega^2$        & $1.44\times10^{-53}\,{\rm m^2}$ & audited bound \\
EHT M87*      & chromatic shadow shift    & $\Lambda\omega^2$        & $2.6\times10^{-9}\,{\rm m^2}$   & validated sensitivity \\
GW (LIGO/LISA)& propagation delay         & $\Lambda\omega^2 K(z)$   & $3.2\times10^{-10}\,{\rm m^2}$  & forecast \\
\bottomrule
\end{tabular}
\end{table}

\begin{remark}[Status definitions --- read carefully]
\textbf{Audited bound (Fermi-LAT):} derived from a published, fitted
constraint on Lorentz-invariance-violating time-of-flight dispersion
\citep{Abdo2009}, with the $\Lambda$-model normalization independently
verified in Section~\ref{sec:fermi}. This is a genuine observational
constraint.

\textbf{Validated sensitivity ceiling (EHT):} the value of $\Lambda$
corresponding to a $10\%$ shadow-size tolerance model calibrated to
published EHT angular resolution at 230\,GHz (Paper 2,
Section~5). This is \emph{not} a fit to actual M87* shadow data; it is
the sensitivity that such a fit \emph{could in principle achieve} under
stated assumptions. A genuine EHT constraint requires per-band
sensitivity, opacity, and imaging-systematics modeling not attempted
here.

\textbf{Forecast (GW):} an illustrative projection of the sensitivity
achievable from a hypothetical multi-band gravitational-wave detection
with electromagnetic counterpart, using the FLRW propagation formula
$\Delta T(\omega,z)=-(\Lambda/2c^3)\,\omega^2\,K(z)$ derived in
Paper 2's FLRW extension. This is \emph{not} a fitted LIGO/Virgo/KAGRA
constraint; no such analysis has been performed on real GW170817 or
other event data.
\end{remark}

All three probes share the identical $\Lambda\omega^2$ scaling, which
follows from the same analytical proof (vanishing of $g^{\mu\nu}k_\mu
k_\nu$ at the relevant null trajectory) established independently in
each spacetime (Schwarzschild for Fermi-LAT and GW forecast, Kerr for
EHT).

%% ════════════════════════════════════════════════════════════════════════════
\section{Condensed-Matter Analog Mappings}
\label{sec:analogs}
%% ════════════════════════════════════════════════════════════════════════════

The flat-space, zero-curvature limit of the dispersion relation is
\begin{equation}\label{eq:flat}
  \omega(k) = ck\sqrt{1+\Lambda k^2}.
\end{equation}
Of the three condensed-matter systems considered during this audit, only
one currently has a literature-supported derivation of this exact form.
We report all three honestly in Table~\ref{tab:mappings}, including the
two that do not (yet) survive scrutiny, rather than omitting the
negative results.

\begin{table}[h]
\centering
\caption{Condensed-matter systems examined for the dispersion form
Eq.~\eqref{eq:flat}. Status reflects the audit in this section, not the
convenience of the algebra.}
\label{tab:mappings}
\begin{tabular}{lllc}
\toprule
System & $\Lambda$ mapping & Physical inputs & Status \\
\midrule
BEC (Bogoliubov)     & $\xi^2/4$              & healing length $\xi$ & \textbf{established} \\
Dirac/Weyl material  & $(1-\eta^2)v_F^2/4$    & warping $\eta$, $v_F$ & \textbf{speculative, unsupported} \\
Photonic ($k$-space) & $\beta_4^k/c^2$        & coeff.\ $\beta_4^k$ of $k^4$ in $\omega^2(k)$ & \textbf{not yet measurable} \\
\bottomrule
\end{tabular}
\end{table}

\subsection{Bose--Einstein condensates}

The Bogoliubov dispersion relation for excitations in a weakly
interacting BEC, $\omega(k)=c_sk\sqrt{1+(\xi k/2)^2}$, matches
Eq.~\eqref{eq:flat} exactly with $\Lambda=\xi^2/4$, verified directly by
matching the $k^4$ coefficients of the two forms
($\omega^2=c_s^2k^2+(\hbar^2/4m^2)k^4$ vs.\ $\omega^2=c_s^2k^2(1+\Lambda
k^2)$). This is an analytic identity, not a numerical coincidence.
\texttt{examples/example\_BEC.py} confirms the fitting code correctly
recovers $\Lambda=\xi^2/4$ from noisy synthetic Bragg-spectroscopy-like
data (healing length $\xi=5\times10^{-7}$\,m, $2\%$ noise, recovered to
$5.4\%$) --- this validates the \emph{fitting procedure}; the physical
mapping itself is independently established by the coefficient-matching
argument above, not by the synthetic-data test.

\subsection{Dirac and Weyl materials --- audited and not supported}
\label{sec:dirac_audit}

An earlier draft of this section proposed $\Lambda=(1-\eta^2)v_F^2/4$
for trigonal/hexagonal-warping-corrected Dirac dispersion, citing
\citet{Fu2009}. On inspection, that citation does not support this
formula. \citet{Fu2009} derives a hexagonal warping term that is
\emph{cubic} in momentum,
\begin{equation}
  H = v_F(k_x\sigma_y - k_y\sigma_x) + \tfrac{\lambda_w}{2}(k_+^3+k_-^3)\sigma_z,
\end{equation}
explicitly described by the author as the counterpart of cubic
Dresselhaus spin-orbit coupling. Squaring to obtain the energy
dispersion gives
\begin{equation}
  E(k)^2 = v_F^2k^2 + \lambda_w^2k^6\cos^2(3\theta),
\end{equation}
an \emph{anisotropic $k^6$} correction, not the isotropic $k^4$
correction required by Eq.~\eqref{eq:flat}. No power of the warping
strength or angular averaging converts a $k^6$ term into a $k^4$ term;
these are structurally different corrections to the dispersion relation.
We were unable to locate any literature derivation of an isotropic
$k^4$ (Lambda-model) correction to Dirac dispersion from trigonal
warping, tilt, or any other known mechanism.

\begin{remark}[On the synthetic-data test]
\texttt{examples/example\_dirac.py} generates synthetic ``ARPES-like''
data directly from the formula $\Lambda=(1-\eta^2)v_F^2/4$ and then
fits that same formula to it. Such a test is guaranteed to succeed
regardless of whether the underlying physical mapping is real; it
demonstrates only that the fitting code is unbiased, and must not be
read as evidence for the mapping. We flag this explicitly because an
earlier version of this manuscript did present the resulting agreement
(``$0.67\%$'') as if it were evidence for the physics. It was not.
\end{remark}

We retain $\Lambda=(1-\eta^2)v_F^2/4$ in the released code, clearly
marked as speculative (it raises a runtime warning on use), for
exploratory purposes only. It should not be treated as a prediction of
this framework, and Table~\ref{tab:mappings} reports it as such.

\subsection{Photonic dispersion --- corrected derivation, not yet
connected to a measurable quantity}

An earlier draft of this section claimed
$\Lambda=\beta_4^k/(12c)$, derived from an expansion
$\omega(k)=ck+(\beta_4^k/24)k^4+\dots$. That expansion is incorrect.
Direct Taylor expansion of Eq.~\eqref{eq:flat} gives
\begin{equation}
  \omega(k) = ck + \tfrac{1}{2}\Lambda c\,k^3 + O(k^5),
\end{equation}
i.e.\ the leading correction to $\omega(k)$ itself is \emph{cubic} in
$k$; there is no $k^4$ term in $\omega(k)$ at any order coming from this
model. The quartic structure appears cleanly only in $\omega^2(k)$,
which is exact and closed-form:
\begin{equation}
  \omega(k)^2 = c^2k^2 + c^2\Lambda\,k^4 \quad\Longrightarrow\quad
  \Lambda = \beta_4^k/c^2, \qquad \beta_4^k \equiv c^2\Lambda,
\end{equation}
where $\beta_4^k$ is defined as the coefficient of $k^4$ in
$\omega^2(k)$ (units ${\rm m^4/s^2}$). This corrected formula is used
throughout the released code.

Even with the corrected algebra, $\beta_4^k$ is not yet connected to any
standard, independently measurable photonics quantity. Standard telecom
dispersion coefficients ($\beta_2,\beta_3,\beta_4=d^n\beta/d\omega^n$)
are Taylor coefficients of the propagation constant expanded around a
nonzero pump/operating frequency --- a structurally different
expansion from the near-$k=0$ form used here. We checked the closest
real physical analogue, quartic dispersion in fiber-optical
event-horizon experiments \citep{Webb2014}: even in the idealized
``pure-quartic'' operating point ($\beta_2=\beta_3=0$), the co-moving-
frame frequency near the group-velocity-matched wavenumber behaves as a
quartic \emph{minimum} with no linear term at all,
$\Omega(\Delta k)\propto-\beta_4\Delta k^4/(24\beta_1^5)$ --- a
different mathematical structure from the deformed-light-cone form of
Eq.~\eqref{eq:flat}, which always retains a leading linear branch
$\omega\to ck$. We do not currently have a derivation connecting
measured fiber $\beta_n$ coefficients to $\beta_4^k$ as used here.

\begin{remark}[On the synthetic-data test]
As in Section~\ref{sec:dirac_audit}, \texttt{examples/example\_photonic.py}
fits synthetic data generated from the same formula it tests, which
validates only the fitting code (now correctly recovering the corrected
$\Lambda=\beta_4^k/c^2$), not any physical claim about real photonic
media. We report this explicitly rather than describing it as
experimental validation.
\end{remark}

%% ════════════════════════════════════════════════════════════════════════════
\section{Open-Source Tools}
\label{sec:tools}
%% ════════════════════════════════════════════════════════════════════════════

All code described in this paper is released under the MIT license in
the \texttt{Lambda-model} repository, organized as:

\begin{verbatim}
Lambda-model/
  paper3/
    wave_equation_2D_solver.py        (Section 2)
    paper3_Dn4_test.py                (Section 2.3)
    paper3_grid_convergence.py        (Section 2.4-2.5)
    lambda_experimental_validator.py  (fitting tool, all domains)
  examples/
    example_BEC.py
    example_dirac.py
    example_photonic.py
\end{verbatim}

The validator (\texttt{lambda\_experimental\_validator.py}) accepts
either raw $(k,\omega)$ CSV data or domain-specific physical parameters,
and reports: the fitted $\Lambda$ with standard error, the fit quality
($R^2$), a significance test against $\Lambda=0$, and (for domain
examples) a direct comparison to the theoretical mapping. This is
intended as a turnkey tool for experimentalists to test their own data
against the framework without needing to reimplement the fitting
procedure. The BEC mapping (\texttt{lambda\_from\_BEC}) is presented as
established; the Dirac mapping (\texttt{lambda\_from\_dirac}) raises a
runtime warning identifying it as speculative and unsupported by the
cited literature (Section~\ref{sec:dirac_audit}); the photonic mapping
(\texttt{lambda\_from\_photonic}) uses the corrected $\beta_4^k/c^2$
form but is explicitly documented as not yet connected to a measurable
laboratory quantity. Domain example scripts print an explicit reminder
that synthetic-data agreement validates the fitting code, not the
physical mapping, wherever that mapping is not independently
established.

%% ════════════════════════════════════════════════════════════════════════════
\section{Conclusions}
%% ════════════════════════════════════════════════════════════════════════════

We have consolidated the quartic-dispersion framework of Papers 1--2
with three additions:

\begin{enumerate}
  \item A validated numerical solver and discrete biharmonic operator,
    with confirmed $O(h^2)$ spatial and $O(1/N)$ temporal convergence,
    and demonstrated $\Lambda$-recovery from noisy data to sub-percent
    accuracy.
  \item A corrected Fermi-LAT bound,
    $\Lambda_{\rm Fermi}=1.4421\times10^{-53}\,{\rm m^2}$, tightened by
    $33\%$ from the original Paper 1 value due to a group-velocity
    normalization error.
  \item One analytically established condensed-matter analog mapping
    (BEC, $\Lambda=\xi^2/4$), with the fitting procedure confirmed
    unbiased on synthetic data; and an honest audit of two further
    proposed mappings (Dirac materials, photonic crystals), neither of
    which currently survives scrutiny as a validated physical claim
    (Section~\ref{sec:analogs}). All fitting tools are released as
    open source regardless of a given mapping's status, so
    experimentalists can test any of them directly against real data.
\end{enumerate}

The overall evidentiary hierarchy is now:
\begin{center}
\textbf{theory (proved)} $\to$ \textbf{numerical validation (this
paper)} $\to$ \textbf{audited bound (Fermi-LAT)} $\to$ \textbf{validated
sensitivity ceiling (EHT)} $\to$ \textbf{forecast (GW)} $\to$
\textbf{open experimental test (condensed matter)}.
\end{center}

We emphasize that no single value of $\Lambda$ has been simultaneously
fitted across all probes; the framework remains a falsifiable hypothesis
about a shared mathematical structure, not a confirmed physical law. The
open-source tools released here are intended to enable that falsification
test directly, by domain experts, on real data.

%% ── Bibliography ────────────────────────────────────────────────────────────
\bibliographystyle{plainnat}
\begin{thebibliography}{99}

\bibitem[Abdo et~al.(2009)]{Abdo2009}
Abdo, A.~A., et~al., 2009, \textit{Astrophys.\ J.}, 706, L138.

\bibitem[Agrawal(2019)]{Agrawal2019}
Agrawal, G.~P., 2019, \textit{Nonlinear Fiber Optics}, 6th ed.\ (Academic Press).

\bibitem[Bogoliubov(1947)]{Bogoliubov1947}
Bogoliubov, N.~N., 1947, \textit{J.\ Phys.\ USSR}, 11, 23.

\bibitem[Fu(2009)]{Fu2009}
Fu, L., 2009, \textit{Phys.\ Rev.\ Lett.}, 103, 266801. (Cited in
Section~\ref{sec:dirac_audit} as the source of the hexagonal warping
term; note that paper does not support the $\Lambda=(1-\eta^2)v_F^2/4$
mapping attributed to it in an earlier draft.)

\bibitem[Webb et~al.(2014)]{Webb2014}
Webb, K.~E., et~al., 2014, \textit{Nature Communications}, 5, 4969
(``Nonlinear optics of fibre event horizons'').

\bibitem[Kretski(2026a)]{Kretski2026a}
Kretski, D., 2026, Zenodo, \texttt{10.5281/zenodo.22018715}
(Paper~1, submitted to CQG, ref.\ CQG-117140).

\bibitem[Kretski(2026b)]{Kretski2026b}
Kretski, D., 2026, Zenodo, \texttt{10.5281/zenodo.22051427}
(Paper~2).

\end{thebibliography}

\end{document}
